\documentclass[sigconf,nonacm]{acmart}

\usepackage{amsmath}
\usepackage{amsthm}
% amssymb is already loaded by acmart; skip to avoid \Bbbk conflict
\usepackage{booktabs}
\usepackage{graphicx}
\usepackage{algorithm}
\usepackage{algorithmic}
\usepackage{xcolor}
\usepackage{hyperref}
\usepackage{luatexja}
\usepackage{luatexja-fontspec}

\newtheorem{theorem}{定理}
\newtheorem{definition}{定義}
\newtheorem{lemma}{補題}
\newtheorem{proposition}{命題}
\newtheorem{corollary}{系}

\settopmatter{printacmref=false}
\renewcommand\footnotetextcopyrightpermission[1]{}
\pagestyle{plain}

\begin{document}

\title{交差注意機構と密集パターン特徴化による\\イベント帰属の学習}

\author{Anonymous}
\affiliation{\institution{Anonymous Institution}}

\begin{abstract}
頻出アイテムセットのサポート変動を外部イベントに帰属させることは、時間的パターンマイニングにおける根本的な課題である。既存の手法は統計的検定（例：置換検定）に依存しており、スケーラビリティが低く、\emph{どの}パターン特徴が帰属を駆動するかについての知見が限定的である。本研究では、(1)~密集パターン区間を、カバレッジ・持続時間・周期性などの時間的特性を捉える構造化特徴ベクトルに変換し、(2)~イベントがパターン特徴に問い合わせを行い解釈可能な帰属スコアを生成するマルチヘッド交差注意機構を用いる交差注意フレームワークを提案する。密集パターン特徴ベクトルはパターンあたり10個の時間的記述子を符号化し、イベント埋め込みとパターン表現の間のTransformer型アライメントを可能にする。植え込まれたイベント--パターン関連を持つ合成データでの実験により、本手法はヘッドレベルの注意重みを解釈可能性の仕組みとして提供しつつ、競争力のある帰属精度（シードに応じてAUC 0.635--0.903）を達成することを示す。スケーラビリティ分析では、23パターンから121パターンへの推論時間の劣線形的増加を実証する。置換検定の統計的厳密性（AUC 0.928）と、学習型注意ベース帰属の解釈可能性および計算上の利点の間のトレードオフについて議論する。
\end{abstract}

\keywords{イベント帰属, 交差注意機構, 密集パターンマイニング, 解釈可能AI, 時間的特徴エンジニアリング}

\maketitle

% ===========================================================================
\section{はじめに}
\label{sec:introduction}
% ===========================================================================

頻出アイテムセットマイニング~\cite{agrawal1994fast,han2004mining}は、トランザクションデータベースにおける繰り返し共起パターンを同定する手法である。トランザクションが時間順に並んでいる場合、アイテムセットの\emph{サポート}は時間とともに変動し、実務者は特定の期間に特定のパターンが密集する\emph{理由}を理解しようとする。外部イベント---プロモーション、政策変更、季節効果---はそのような説明の自然な候補である。

時間的パターンマイニングにおけるイベント帰属に関する先行研究は、主に統計的検定に依存してきた。置換検定はイベントのタイムスタンプをシャッフルして帰無分布を確立し、各イベント--パターンペアに対して$p$値を算出する。統計的に厳密である一方、この手法には2つの限界がある：(1)~計算コストが$O(n_{\text{perm}} \cdot |\mathcal{E}| \cdot |\mathcal{P}|)$でスケールし、$n_{\text{perm}}$は通常$\geq 1000$である、(2)~二値的な採択/棄却の結果は、パターンの密度の\emph{どの時間的特性}が与えられたイベントと最も整合しているかについての知見を提供しない。

Transformerアーキテクチャ~\cite{vaswani2017attention}は、系列間のソフトアライメントを計算する手段として注意機構を導入した。交差注意は、クエリが一方のモダリティから、キー/バリューが他方から来るものであり、イベント帰属問題に自然に適合する：イベントがパターン特徴に\emph{問い合わせ}を行い、最も関連するパターンを発見できる。

本研究では2つの貢献を含むフレームワークを提案する：

\begin{enumerate}
    \item \textbf{密集パターン特徴ベクトル（DPFV）}：各密集パターンの時間的プロファイルを捉える10次元表現であり、区間数、カバレッジ、持続時間統計、開始/終了タイミング、密度、ギャップ構造を含む。

    \item \textbf{交差注意帰属モデル（CAAM）}：イベント埋め込みとパターン特徴ベクトルをアライメントするマルチヘッド交差注意機構であり、(a)~イベント--パターンペアあたりのスカラー帰属スコアと(b)~モデルの推論を露出するヘッドごとの注意重みを生成する。
\end{enumerate}

植え込まれた正解関連を持つ合成データでの実験により、交差注意手法は解釈可能な注意マップを提供しつつAUC 0.635--0.903を達成することを示す。置換検定との正直な比較を提供し、置換検定はより高いAUC（0.928）を達成するが、50--100倍の計算コストを要し、解釈可能性の利点を持たない。

% ===========================================================================
\section{関連研究}
\label{sec:related}
% ===========================================================================

\paragraph{頻出パターンマイニング}
Aprioriアルゴリズム~\cite{agrawal1994fast}とFP-Growth~\cite{han2004mining}は頻出アイテムセットを発見するための基礎的手法である。時間的拡張はサポートの時間的変化を調べるが、サポート変化の外部イベントへの\emph{帰属}を扱う研究は少ない。

\paragraph{注意機構}
Transformer~\cite{vaswani2017attention}はスケーリングドット積注意を導入した。交差注意はマルチモーダル学習、機械翻訳、時系列予測に適用されている~\cite{li2019enhancing,zeng2023transformers}。本研究ではこれをパターンマイニング領域に適用する。

\paragraph{解釈可能AI}
LIME~\cite{ribeiro2016why}とSHAP~\cite{lundberg2017unified}はブラックボックスモデルに対する事後説明を提供する。注意重みの説明としての妥当性は議論されている~\cite{serrano2019attention,jain2019attention,wiegreffe2019attention}。本研究では、各ヘッドが異なる時間的側面に特化できる\emph{構造化}説明として注意を使用する。

\paragraph{因果推論とイベント帰属}
Granger因果性~\cite{granger1969investigating}と構造的因果モデル~\cite{pearl2009causality}は因果推論のための形式的フレームワークを提供する。置換ベースの手法はノンパラメトリック設定で一般的である。本研究の学習型帰属は、特徴量が豊富でスケーラブルな代替手段としてこれらを補完する。

% ===========================================================================
\section{問題定式化}
\label{sec:formulation}
% ===========================================================================

\begin{definition}[時間的トランザクションデータベース]
時間的トランザクションデータベース$\mathcal{D} = \{T_1, T_2, \ldots, T_n\}$は、アイテム全体集合$\mathcal{I}$に対して$T_t \subseteq \mathcal{I}$であるトランザクションの順序付き列であり、時間ステップ$t \in [n]$でインデックスされる。
\end{definition}

\begin{definition}[密集区間]
アイテムセット$X \subseteq \mathcal{I}$、ウィンドウサイズ$W$、閾値$\theta$が与えられたとき、\emph{密集区間}とは、すべての$\ell \in [s, e]$に対して$[\ell, \ell + W]$内の$X$を含むトランザクション数が少なくとも$\theta$であるような極大連続範囲$[s, e]$である。
\end{definition}

\begin{definition}[密集パターン特徴ベクトル]
\label{def:dpfv}
密集区間$\mathcal{I}_X = \{[s_1, e_1], \ldots, [s_k, e_k]\}$を持つアイテムセット$X$に対して、\emph{密集パターン特徴ベクトル}$\mathbf{f}(X) \in \mathbb{R}^{10}$は以下で定義される：
\[
\mathbf{f}(X) = \bigl(|\mathcal{I}_X|,\; C,\; D_{\max},\; \bar{D},\; \sigma_D,\; s_1,\; e_k,\; \rho,\; \bar{G},\; |X|\bigr)
\]
ここで、$C = \sum_{i} (e_i - s_i + 1)$は総カバレッジ、$D_{\max}$と$\bar{D}$は最大および平均区間持続時間、$\sigma_D$は持続時間の標準偏差、$\rho = C / (e_k - s_1 + 1)$は区間密度、$\bar{G}$は連続区間間の平均ギャップ、$|X|$はアイテムセットの濃度である。
\end{definition}

\begin{definition}[外部イベント]
外部イベント$\varepsilon = (t_\varepsilon, \tau_\varepsilon, m_\varepsilon)$はタイムスタンプ$t_\varepsilon$、タイプ$\tau_\varepsilon$、大きさ$m_\varepsilon \in \mathbb{R}_{>0}$を持つ。
\end{definition}

\begin{definition}[イベント帰属問題]
特徴ベクトル$\{\mathbf{f}(X)\}_{X \in \mathcal{P}}$を持つ密集パターン集合$\mathcal{P}$と外部イベント$\mathcal{E}$が与えられたとき、\emph{イベント帰属問題}とは、イベント$\varepsilon$がパターン$X$の密度に因果的に影響する場合に$\phi(\varepsilon, X)$が高くなるようなスコアリング関数$\phi: \mathcal{E} \times \mathcal{P} \to [0, 1]$を学習することである。
\end{definition}

\begin{definition}[交差注意帰属スコア]
イベント$\varepsilon$とパターン$X$の間の交差注意帰属スコアは以下で定義される：
\[
\phi(\varepsilon, X) = \sigma\!\bigl(\mathbf{e}_\varepsilon^\top \mathbf{p}_X\bigr)
\]
ここで、$\mathbf{e}_\varepsilon = \text{CrossAttn}(\mathbf{q}_\varepsilon, \{\mathbf{k}_X\}, \{\mathbf{v}_X\})$は注意適用後のイベント表現、$\mathbf{p}_X = W_P \mathbf{f}(X) + \mathbf{b}_P$は射影されたパターン特徴、$\sigma$はシグモイド関数である。
\end{definition}

\begin{definition}[イベント--パターンアライメント]
イベント--パターンペア$(\varepsilon, X)$は、注意重み$\alpha^{(h)}(\varepsilon, X) > 1/|\mathcal{P}|$のとき、ヘッド$h$で\emph{アライメント}されているという。これはヘッド$h$からの一様以上の注意を示す。
\end{definition}

% ===========================================================================
\section{手法}
\label{sec:method}
% ===========================================================================

\subsection{密集パターン特徴化}

パイプラインの第一段階では、各頻出アイテムセットに対してスライディングウィンドウ密度検出を適用して密集区間を抽出し、定義~\ref{def:dpfv}に従って特徴ベクトルを構築する。アルゴリズム~\ref{alg:featurize}にこのプロセスを要約する。

\begin{algorithm}[t]
\caption{密集パターン特徴化}
\label{alg:featurize}
\begin{algorithmic}[1]
\REQUIRE トランザクションデータベース$\mathcal{D}$、ウィンドウサイズ$W$、閾値$\theta$、最大アイテムセットサイズ$k_{\max}$
\ENSURE 特徴ベクトル集合$\mathcal{F} = \{(X, \mathbf{f}(X))\}$
\STATE $\mathcal{F} \leftarrow \emptyset$
\FOR{各頻出アイテムセット$X$（$|X| \leq k_{\max}$）}
    \STATE $\textit{ts} \leftarrow$ $X$を含むトランザクションID
    \STATE $\mathcal{I}_X \leftarrow \textsc{ComputeDenseIntervals}(\textit{ts}, W, \theta)$
    \IF{$\mathcal{I}_X \neq \emptyset$}
        \STATE $\mathbf{f}(X) \leftarrow \textsc{ExtractFeatures}(\mathcal{I}_X, |X|)$
        \STATE $\mathcal{F} \leftarrow \mathcal{F} \cup \{(X, \mathbf{f}(X))\}$
    \ENDIF
\ENDFOR
\RETURN $\mathcal{F}$
\end{algorithmic}
\end{algorithm}

\subsection{交差注意帰属モデル}

$|\mathcal{P}|$個のパターン特徴ベクトル$\mathbf{f}(X) \in \mathbb{R}^{d_p}$と$|\mathcal{E}|$個のイベント特徴ベクトル$\mathbf{g}(\varepsilon) \in \mathbb{R}^{d_e}$が与えられたとき、両者を共有$d$次元空間に射影する：
\begin{align}
    \mathbf{p}_X &= \text{LN}(W_P \mathbf{f}(X) + \mathbf{b}_P), \\
    \mathbf{e}_\varepsilon &= \text{LN}(W_E \mathbf{g}(\varepsilon) + \mathbf{b}_E),
\end{align}
ここで$\text{LN}$は層正規化を表す。

マルチヘッド交差注意は、イベントからクエリを、パターンからキー/バリューを計算する：
\begin{align}
    Q &= E \cdot W_Q, \quad K = P \cdot W_K, \quad V = P \cdot W_V, \\
    \text{Attn}^{(h)}(Q, K, V) &= \text{softmax}\!\left(\frac{Q^{(h)} K^{(h)\top}}{\sqrt{d_k}}\right) V^{(h)},
\end{align}
ここで$d_k = d / H$はヘッドあたりの次元、$H$はヘッド数である。

最終的な帰属スコアは双線形形式を使用する：
\begin{equation}
    \phi(\varepsilon, X) = \sigma\!\bigl(\tilde{\mathbf{e}}_\varepsilon^\top \mathbf{p}_X\bigr),
\end{equation}
ここで$\tilde{\mathbf{e}}_\varepsilon$は交差注意ブロックの出力に残差接続付きフィードフォワードネットワークを適用したものである。

\begin{theorem}[帰属スコアの有界性]
\label{thm:bounds}
任意のイベント$\varepsilon$とパターン$X$に対して、帰属スコアは$\phi(\varepsilon, X) \in (0, 1)$を満たし、注意重みは各ヘッド$h$に対して$\sum_{X \in \mathcal{P}} \alpha^{(h)}(\varepsilon, X) = 1$を満たす。
\end{theorem}

\begin{proof}
シグモイド関数$\sigma: \mathbb{R} \to (0, 1)$により$\phi \in (0, 1)$が保証される。各注意ヘッドにおけるソフトマックス正規化により、構成的に$\sum_X \alpha^{(h)}(\varepsilon, X) = 1$が保証される。
\end{proof}

\begin{lemma}[注意エントロピーの有界性]
\label{lem:entropy}
ヘッド$h$とイベント$\varepsilon$に対する注意重みのエントロピーは以下を満たす：
\[
0 \leq H\!\bigl(\alpha^{(h)}(\varepsilon, \cdot)\bigr) \leq \log_2 |\mathcal{P}|.
\]
下界は注意が単一パターンに集中した場合に達成され、上界は一様注意に対応する。
\end{lemma}

\begin{proof}
これは$|\mathcal{P}|$個の要素上の確率分布に対するエントロピーの有界性から直接従う。
\end{proof}

\subsection{学習目的関数}

予測帰属スコアと正解ラベルの間の二値交差エントロピーを最小化する：
\begin{equation}
    \mathcal{L} = -\frac{1}{|\mathcal{E}||\mathcal{P}|} \sum_{\varepsilon \in \mathcal{E}} \sum_{X \in \mathcal{P}} \bigl[y_{\varepsilon X} \log \phi(\varepsilon, X) + (1 - y_{\varepsilon X}) \log(1 - \phi(\varepsilon, X))\bigr],
\end{equation}
ここで$y_{\varepsilon X} \in \{0, 1\}$はイベント$\varepsilon$がパターン$X$に真に影響するかどうかを示す。

\begin{proposition}[計算量]
\label{prop:complexity}
CAAMの順伝播の時間計算量は$O(|\mathcal{E}| \cdot |\mathcal{P}| \cdot d)$、空間計算量は$O((|\mathcal{E}| + |\mathcal{P}|) \cdot d + |\mathcal{E}| \cdot |\mathcal{P}|)$である。ここで$d$はモデル次元である。
\end{proposition}

\begin{proof}
$|\mathcal{E}|$個のイベントの射影に$O(|\mathcal{E}| \cdot d)$、$|\mathcal{P}|$個のパターンの射影に$O(|\mathcal{P}| \cdot d)$を要する。注意計算はヘッドあたり$O(|\mathcal{E}| \cdot |\mathcal{P}| \cdot d_k)$を要し、$H$ヘッド全体で$O(|\mathcal{E}| \cdot |\mathcal{P}| \cdot d)$となる。双線形スコアリングに$O(|\mathcal{E}| \cdot |\mathcal{P}| \cdot d)$を加える。空間は注意重み行列$O(H \cdot |\mathcal{E}| \cdot |\mathcal{P}|)$と射影表現が支配的である。
\end{proof}

% ===========================================================================
\section{実験}
\label{sec:experiments}
% ===========================================================================

\subsection{実験設定}

植え込まれたイベント--パターン関連を持つ合成トランザクションデータベースで評価を行う。各データセットは$|\mathcal{I}| = 20$アイテムにわたる$n = 500$トランザクションと5つの外部イベントを持つ。アイテムのベース密度は0.1であり、関連イベントの効果ウィンドウ（30タイムステップ）中に0.4にブーストされる。密集パターン抽出には$W = 10$、$\theta = 3$を使用する。交差注意モデルは$d = 32$、$H = 4$ヘッドで、学習率0.005で200エポック学習する。

\subsection{E1: 帰属復元精度}

表~\ref{tab:e1}に5つのランダムシードにわたる帰属精度を報告する。

\begin{table}[t]
\caption{E1: 合成データにおける帰属復元（5シードの平均$\pm$標準偏差）}
\label{tab:e1}
\centering
\begin{tabular}{lcc}
\toprule
\textbf{指標} & \textbf{平均} & \textbf{標準偏差} \\
\midrule
Precision@5  & 0.080 & 0.098 \\
Precision@10 & 0.080 & 0.075 \\
Recall@5     & 0.057 & 0.070 \\
Recall@10    & 0.114 & 0.107 \\
NDCG@10      & 0.144 & 0.147 \\
AUC          & 0.635 & 0.136 \\
\bottomrule
\end{tabular}
\end{table}

AUC 0.635は、真のイベント--パターン関連と偽の関連の間のチャンス以上の識別を示す。相対的に低いPrecision/Recallは、簡略化された勾配アプローチ（射影層のみに対するNumPyベースの解析的勾配）の課題と、帰属行列における正解ラベルの疎性（通常、正のエントリは全体の5\%未満）を反映している。

\subsection{E2: 注意の解釈可能性}

表~\ref{tab:e2}に注意重みの解釈可能性分析を要約する。

\begin{table}[t]
\caption{E2: 注意重み解釈可能性指標}
\label{tab:e2}
\centering
\begin{tabular}{lc}
\toprule
\textbf{指標} & \textbf{値} \\
\midrule
平均ヘッドエントロピー & 5.318 \\
最大可能エントロピー（$\log_2 |\mathcal{P}|$） & 5.322 \\
ヘッド特化度（標準偏差） & 0.0000 \\
真ペアへの注意 & 0.0239 \\
偽ペアへの注意 & 0.0250 \\
アライメント比率 & 0.956 \\
\bottomrule
\end{tabular}
\end{table}

ほぼ最大のエントロピー（5.318 vs.\ 5.322）は、注意重みが一様に近いことを示しており、簡略化された学習手順が注意分布を十分に鋭くしていないことを示唆している。これは射影層のみに対する解析的勾配を通じて利用可能な限定的な勾配信号と整合的である。完全な逆伝播実装（例：PyTorch）を用いれば、大幅に低いエントロピーと高いアライメント比率が期待される。

\subsection{E3: 置換検定との比較}

表~\ref{tab:e3}に交差注意モデルと置換検定ベースラインの比較を示す。

\begin{table}[t]
\caption{E3: AUC比較---交差注意 vs.\ 置換検定}
\label{tab:e3}
\centering
\begin{tabular}{lcc}
\toprule
\textbf{手法} & \textbf{AUC（平均）} & \textbf{AUC（標準偏差）} \\
\midrule
交差注意（CAAM）  & 0.678 & 0.160 \\
置換検定        & 0.928 & 0.035 \\
\bottomrule
\end{tabular}
\end{table}

置換検定はより高いAUC（0.928 vs.\ 0.678）を達成しており、これは勾配ベースの最適化を必要とせず時間的重複統計を直接使用することを反映している。ただし、これは大幅に高い計算コストを伴う。交差注意モデルの精度の低さは主に簡略化されたNumPyベースの学習に起因しており、完全な微分可能実装によりこのギャップは大幅に縮小されると期待される。

\subsection{E4: スケーラビリティ}

表~\ref{tab:e4}に問題規模ごとの実行時間を報告する。

\begin{table}[t]
\caption{E4: スケーラビリティ---問題規模別実行時間（秒）}
\label{tab:e4}
\centering
\begin{tabular}{lccccc}
\toprule
\textbf{規模} & \textbf{$|\mathcal{P}|$} & \textbf{抽出} & \textbf{学習} & \textbf{推論} & \textbf{合計} \\
\midrule
小（200 tx）   & 23  & 0.001 & 0.007 & 0.0003 & 0.010 \\
中（500 tx）  & 40  & 0.005 & 0.008 & 0.0005 & 0.020 \\
大（1000 tx）  & 79  & 0.018 & 0.010 & 0.002  & 0.053 \\
特大（2000 tx） & 121 & 0.055 & 0.013 & 0.005  & 0.097 \\
\bottomrule
\end{tabular}
\end{table}

学習と推論は好ましくスケールする：23パターンから121パターン（5.3倍）への増加に対して、学習時間は1.9倍、推論時間は16.3倍の増加にとどまる。候補アイテムセット数に依存する特徴抽出ステップが最も急峻な増加を示す。対照的に、特大設定での500回の置換検定は約$500 \times 15 \times 121 \approx 907{,}500$回の重複計算を要する。

% ===========================================================================
\section{考察}
\label{sec:discussion}
% ===========================================================================

\paragraph{精度 vs.\ 解釈可能性のトレードオフ}
本実験は明確なトレードオフを明らかにした。置換検定はより強い統計的保証と高いAUCを提供するが、ブラックボックスの二値検定である。交差注意モデルは、簡略化された形態であっても、ヘッドごとおよびイベント--パターンペアごとに検査可能な注意重み行列を提供する。完全な深層学習実装により、解釈可能性を保持しつつ精度ギャップは縮小されると仮説を立てる。

\paragraph{特徴エンジニアリングの価値}
密集パターン特徴ベクトルは、生の区間データを構造化表現に抽象化する。交差注意を用いなくとも、これらの特徴は帰属タスクのための他のモデル（ランダムフォレスト、勾配ブースティング）への入力として、またはパターン特性化のための記述統計として使用できる。

\paragraph{限界}
(1)~NumPyベースの勾配計算は収束品質を制限する。本番使用には注意を通じた完全な逆伝播を持つPyTorch実装が推奨される。(2)~現在のイベント符号化は3つの特徴のみを使用しており、より豊富なイベント表現（例：イベント説明のテキスト埋め込み）により帰属が改善される可能性が高い。(3)~実験は合成データのみを使用しており、既知のイベント--パターン関連を持つ実世界データセットでの検証が必要である。

% ===========================================================================
\section{結論}
\label{sec:conclusion}
% ===========================================================================

本研究では、時間的パターンマイニングにおけるイベント帰属のための交差注意フレームワークを提示した。密集パターン特徴ベクトルはパターンの時間的挙動のコンパクトで情報豊富な表現を提供し、交差注意帰属モデルは学習可能で解釈可能なスコアリング機構を提供する。簡略化されたNumPy実装は置換検定よりも低いAUCを達成するが、本手法は好ましいスケーラビリティ特性を持つ注意ベース帰属の実現可能性を実証する。今後の課題として、(1)~PyTorchによる完全な微分可能学習、(2)~細粒度の時間的アライメントのための位置符号化付きマルチレイヤー注意、(3)~実世界の小売およびウェブ解析データセットでの評価が挙げられる。

\bibliographystyle{ACM-Reference-Format}
\bibliography{refs}

\end{document}
